\documentclass{situationist-preprint}

\title{Factivity and Self-Reference}
\msid{LOG-10-b}
\msfamily{LOG-10}
\msdiscipline{Logic and Philosophy of Logic}
\mstype{Companion manuscript}
\mskeywords{paradox of the knower; factivity; self-reference; blindspots; Löb's theorem; reflection}

\begin{document}
\maketitle

\begin{abstract}
LOG-10-a argues that the paradox of the knower shows that knowledge, treated as a predicate of sentences that speak about it, cannot satisfy both unrestricted factivity and closure, and that the premise to restrict is the knower's knowledge of its own factivity. We develop the role of factivity itself. Factivity alone, together with the diagonal sentence, implies that the knower sentence is true and not known. What the paradox exhibits is therefore not an incoherence in factivity but a true sentence that a particular knower cannot know, while others can. The knower sentence is a self-referential blindspot, and the restriction on known factivity is the condition under which the knower does not violate it.
\end{abstract}

\section{What factivity proves}

Let $K_a$ be a predicate of sentences, read \enquote{is known by $a$}, in a language containing Robinson arithmetic, and let $S_a$ be a sentence such that $S_a \leftrightarrow \neg K_a\ulcorner S_a\urcorner$ is provable, as in equation~(1) of LOG-10-a, following \citet{KaplanMontague1960}. We write (i)--(iv) for the schemata of Definition~2.1 of LOG-10-a, relativized to $a$.

\begin{proposition}\label{p:true}
Factivity for the single sentence $S_a$, that is, $K_a\ulcorner S_a\urcorner \rightarrow S_a$, together with the diagonal biconditional, logically implies $S_a \wedge \neg K_a\ulcorner S_a\urcorner$.
\end{proposition}

\begin{proof}
Suppose $K_a\ulcorner S_a\urcorner$. By factivity, $S_a$, and by the biconditional, $\neg K_a\ulcorner S_a\urcorner$. So $\neg K_a\ulcorner S_a\urcorner$, and by the biconditional, $S_a$.
\end{proof}

Proposition~\ref{p:true} is the first half of the proof of Theorem~3.1 of LOG-10-a, isolated. It uses no closure, no known arithmetic and no known factivity, and it derives no contradiction. It says that if $a$'s knowledge is factive, then $S_a$ is a truth that $a$ does not know. Factivity, in the presence of self-reference, is not threatened; it is productive. It yields a truth about $a$, and identifies it.

\begin{corollary}\label{c:blind}
Given factivity, closure and known arithmetic for $a$, $a$ knows the instance $K_a\ulcorner S_a\urcorner \rightarrow S_a$ only if the schemata are inconsistent. A second knower $b$, whose knowledge is factive and who knows that $a$'s knowledge is factive for $S_a$ and knows the diagonal biconditional, knows $S_a$, if $b$'s knowledge is closed.
\end{corollary}

\begin{proof}
The first claim is Theorem~3.1 of LOG-10-a. For the second, $b$ knows the two premises of Proposition~\ref{p:true}, and its conclusion follows from them by logic; by closure for $b$, $b$ knows $S_a \wedge \neg K_a\ulcorner S_a\urcorner$ and so $S_a$. No contradiction arises, because $S_a$ says that $a$, not $b$, does not know it.
\end{proof}

Corollary~\ref{c:blind} is the epistemic counterpart of a familiar fact about formal theories. A sound theory $U$ does not prove its Gödel sentence, while we, reasoning about $U$ from the premise that $U$ is sound, know that the sentence is true. The restriction of Proposition~3.2 of LOG-10-a is exactly this: $U$ satisfies factivity, closure and known arithmetic, and fails known factivity at the instances that would let it reproduce, about itself, the reasoning that others can carry out about it. Löb's theorem \citep{Lob1955} fixes which instances these are.

\section{Self-referential blindspots}

\citet{Sorensen1988} calls a proposition a blindspot for a subject if it is consistent but cannot be known by that subject. The knower sentence, relativized to $a$, is a blindspot for $a$: by Proposition~\ref{p:true} it is true if $a$'s knowledge is factive, and by Corollary~\ref{c:blind} it can be known by others. The surprise announcement of LOG-11 is a blindspot for the students on a given morning, and the Fitch conjunction $p \wedge \neg Kp$ of LOG-19 is a blindspot for every knower, since $K$ there quantifies over all knowers and times. The knower sentence differs from both in that it is generated by self-reference, and so requires a predicate; but its epistemic profile is the same. In each case, knowledge by the subject would falsify what is known, and factivity forbids it.

This yields a further reason for the choice of LOG-10-a to restrict known factivity rather than factivity or closure. A knower who cannot know its own blindspots loses nothing by failing to know the instance of factivity that would let it derive one. The failure is not a gap in its self-knowledge that better reflection could close: by Corollary~\ref{c:blind}, any knower who could close it would know a blindspot of its own, which is impossible.

\begin{objection}[Impersonal knowledge]
If $K$ means \enquote{is known by someone at some time}, there is no second knower. Then Proposition~\ref{p:true} shows that $S$ is true, and we know it, which is the paradox again.
\end{objection}

\begin{reply}
Proposition~\ref{p:true} shows that $S$ follows from factivity for $S$. If $K$ is impersonal, no one can know that instance of factivity, by Theorem~3.1 of LOG-10-a; so no one can know $S$ by this route, ourselves included. What we can know is the conditional: if knowledge is factive for $S$, then $S$ is true and unknown. On the impersonal reading, the instance of factivity for $S$ is a truth that no one knows, and the knower paradox yields an unknowable truth by the route of self-reference, as Fitch's argument yields one without it. This agrees with LOG-19 that not every truth is knowable.
\end{reply}

\section{Conclusion}

Factivity, applied to the knower sentence, proves that the sentence is true and unknown to the knower it concerns. The contradiction of the paradox arises only when the knower is credited with knowing that instance of its own factivity. The knower sentence is a blindspot, knowable by others but not by its subject, and the restriction on known factivity defended in LOG-10-a is what it takes for a knower not to know its own blindspots.

\begin{thebibliography}{99}\small
\bibitem[Kaplan and Montague(1960)]{KaplanMontague1960} Kaplan, D. and Montague, R. (1960). A paradox regained. \emph{Notre Dame Journal of Formal Logic}, 1(3), 79--90.
\bibitem[Löb(1955)]{Lob1955} Löb, M.~H. (1955). Solution of a problem of Leon Henkin. \emph{The Journal of Symbolic Logic}, 20(2), 115--118.
\bibitem[Sorensen(1988)]{Sorensen1988} Sorensen, R.~A. (1988). \emph{Blindspots}. Clarendon Press.
\end{thebibliography}

\end{document}
